Door haar kortgesloten
denkvermogen schoten

weer de bliksemschichten,
die haar deden zwichten.

De intense hoofdpijn
ontsteeg haar bewustzijn.

De lichtgevoeligheid
en de misselijkheid

verstoorden haar visie.
Ze gleed in regressie.

Plotinus kon baden
in zijn Enneaden

in de vlucht van zijn ziel
nadat het even viel,

zoals Emmy daalde
en zich daar ophaalde.

Haar herinneringen
in de duizelingen

helderden steeds meer op.
De tijd stond op zijn kop,

toen haar brein oplaaide.
Een raddraaier draaide

het rad van samsara
vele decennia

tot ze centraliteit
ervoer als aankomsttijd.

In warmte dreef ze rond,
in water kalm gegrond.

Geen nijd, geen vrees, geen lust.
Slechts tijd, slechts vlees, slechts rust

brachten Emmy terug,
naar waar ze vliegensvlug

een blik op mocht werpen
om die aan te scherpen.

De gevoelsenergie
trok met zijn frequentie

een vorig leven aan.
Daar bleef de tijd stil staan.

De wet van aantrekking
heeft immers betrekking

op voornamelijk lust.
Die had ze uitgeblust

in haar eerste leven.
Genot opgegeven.

In huis was het rustig.
Maria gleed lustig

zonder ondersteuning
over een trapleuning

met links en rechts een been.
Moeder strafte meteen

en schreeuwde volmondig:
“Dit genot is zondig”.

Niemand, zelfs niet eens God,
zag ooit Maria’s grot

(buiten Bernadette).
Dat was etiquette.

Als gifdampen stonken,
uit diepe spelonken

verborgen voor het licht,
plamuurde zij ze dicht.

Ze had als lichtwezen
haar rol niet gelezen.

Maria daardoor had
een onbeschreven blad,

fris, vers, dun en van goud.
Nog groen in plaats van oud.

Ze leefde als Barbie.
Schijnbare harmonie

en neppe perfectie
gaf de incarnatie

weinig zin en diepgang.
Het duurde ook niet lang

of de sprookjeswereld
werd drastisch ontregeld.

Het netjes ingekleurd
werd turbulent besmeurd.

Haar ma pleegde zelfmoord,
na nooit te zijn gehoord,

door pa die priester was
en niemand echt genas

van het aardse lijden.
De schande bestrijden

was zijn voornaamste doel.
Verder bleef hij vrij koel

onder het gebeuren.
Hij liet zich meesleuren

in zijn echte waarde;
reputatieschade.

Dus kwam er een wit graf,
dat hem zijn aanzien gaf,

maar ook doodsbeenderen
onder zijn eenderen.

Het irriteerde haar.
Hij was een huichelaar

ingevolge Jezus,
zo schreef ooit Matteüs.

Maria werd slechts non,
omdat het helpen kon

va’s eer te herstellen.
Het oordeel te vellen

door de geestelijkheid:
dat vader geen verwijt

of schuld, maar een ramp trof.
Het was geen gespreksstof

meer voor de gemeenschap,
wat betreft vaderschap.

Moeder was bezeten.
Pa werd niets verweten.

Ze zag dit leven aan
als een begrensd bestaan

dat uit volheid begon,
maar Emmy overwon

door de smorende daad
van vechten tegen kwaad.

Derhalve koos Emmy
te gaan voor transgressie

in een volgend leven.
Zich daar vol te geven.

“Sorry, daar zit ik niet”,
zei de veinzende griet.

Haar medezuster vroeg
hoe ze de blik verdroeg

van de knappe tuinman.
Daar lustte ze meer van.

Hij wist te verhitten.
“Je kunt niet vastzitten.

Wie zijn geest begeleid,
wandelt immers altijd.”,

zei de meer geleerde.
Maria negeerde

de genitaliën.
De evangeliën

ontnamen de Barbies
hun spleet met bloedverlies

en de Kents de lullen,
die de spleetjes vullen.

Het bestaan uitgeblust.
Maria had haar lust

niet zozeer afgestopt.
Ze was slechts afgestompt

en reageerde kil.
De zuster zei: “Ik wil

dat je er voor me bent.”
Maar Maria ontkent;

“Ik ben toch dicht bij je?"
Ze speelde het klunsje,

maar wist vast wel beter.
Werd een gesprek heter,

dan blokkeerde ze vaak.
Dat was een schrale zaak.

Dus vingerde ze zich
in haar cel gluiperig
om nog iets te voelen.
Haar nood weg te spoelen.

In de armste oorden
klonken vaak haar woorden

en de collectebus.
Namens Benedictus

vroeg ze om aalmoezen.
Ze gebruikte smoezen

om te verdoezelen
wie te vertroetelen

met de vrome giften.
De kerk kon meeliften

op geld- of hongersnood.
Menigeen nam aanstoot

aan deze praktijken
van zichzelf verrijken.

Zo stond eens Maria
als steelse paria

voor de deuropening
van een zeer arm gezin.

“Een gift voor de armen?”
In zijn grote armen

tilde het gezinshoofd
wat ze al had geroofd

ver boven zichzelf uit
en verkondigde luid:

“Dank U goedgeefse God
dat U zich nu ons lot

eindelijk eens aantrekt.”
Zijn kas werd toen gespekt

met alles uit de bus.
Haar voorhoofd kreeg een kus,

waarna vlug de proleet
zijn voordeur hard dicht smeet.

Maria leek volleerd
en geplastificeerd.

Als Barbie verworden
om bemind te worden

en dan op te merken
dat het zou beperken

als ze geen sex meer had.
Haar nep dempte haar gat,

de diepte in een mens.
Eigenheid was haar wens,

maar zocht een schouderklop.
Het leek er misschien op,

maar het was niet gelukt.
Ze verwerd een construct

vast in een barbiepop.
Carl Rogers wees erop:

“Je authenticiteit
raak je pertinent kwijt

als je slechts kan leven
door steeds toe te geven

aan wat men van je eist.
Je ware zelf afwijst.”

Emmy zag dit al lang.
Iedereen leek zo bang

om echt zichzelf te zijn.
Gevangen in hun pijn

werden ze als zombies,
hersenloze junkies

zoekend naar erkenning
door de samenleving.

Zelfs in een hippie oord
deed men zoals het hoort.

Ze kochten tupperware,
immers dan was men er.

Waar belandde Emmy
met haar retrospectie?
